\chapter*{λ-calculus / type theory to Lean dictionary}
\addcontentsline{toc}{chapter}{λ-calculus / type theory to Lean dictionary}
\setcounter{section}{0}\label{sec:root:lambda-calculus-dictionary:1}

A direct dictionary connecting the formal λ-calculus/type-theory notation used in a few ``Mathematical reading'' boxes throughout this book back to the Lean syntax it corresponds to. This page is a lookup table, not something to read start to finish. \hyperref[sec:02-terminology-and-coc:01-terminology:1]{Chapter 2, Section 1} and \hyperref[sec:02-terminology-and-coc:02-pi-sigma-and-coc:1]{Chapter 2, Section 2}, \hyperref[sec:04-propositions-and-proofs:02-logic-recap:1]{Chapter 4, Section 2}, and \hyperref[sec:06-rigor-check:03-typing-rules-and-safety:1]{Chapter 6, Section 3} are where each row is actually built up and explained.

\begin{longtable}[]{@{}
  >{\raggedright\arraybackslash}p{(\linewidth - 4\tabcolsep) * \real{0.3333}}
  >{\raggedright\arraybackslash}p{(\linewidth - 4\tabcolsep) * \real{0.3333}}
  >{\raggedright\arraybackslash}p{(\linewidth - 4\tabcolsep) * \real{0.3333}}@{}}
\toprule\noalign{}
\begin{minipage}[b]{\linewidth}\raggedright
λ-calculus / type theory
\end{minipage} & \begin{minipage}[b]{\linewidth}\raggedright
Lean syntax
\end{minipage} & \begin{minipage}[b]{\linewidth}\raggedright
Where in this book
\end{minipage} \\
\midrule\noalign{}
\endhead
\bottomrule\noalign{}
\endlastfoot
variable \(x\) & \texttt{x} & Chapter 1 \\
abstraction \(\lambda x.\, t\) & \texttt{fun x => t} & Chapter 1 \\
application \(t_1\, t_2\) & \texttt{t1 t2} & Chapter 1 \\
α-conversion & (silent, automatic) & --- \\
β-reduction \((\lambda x.t)s \to t[x:=s]\) & \texttt{rfl}, \texttt{\#eval} reduction & Chapter 2, Section 1, Chapter 6, Section 4 (definitional equality) \\
type judgment \(\Gamma \vdash t : \tau\) & \texttt{\#check t} & Chapter 1 \\
simple function type \(\tau_1 \to \tau_2\) & \texttt{τ1 → τ2} & Chapter 1, Chapter 6, Section 3 \\
universe hierarchy \(\mathtt{Type}\,i\) & \texttt{Type}, \texttt{Type 1}, \ldots{} & Chapter 6, Sections 2--3 \\
Π-type \(\prod_{x:A} B(x)\) & \texttt{(x : A) → B x}, \texttt{∀ x : A, B x} & Chapter 1, Sections 3, 5, Chapter 4 \\
Σ-type \(\sum_{x:A} B(x)\) & \texttt{structure} (a general dependent-pair type, extractable), of which \texttt{Sigma}/the anonymous-constructor \texttt{⟨\_\allowbreak{}, \_\allowbreak{}⟩} pairing is a direct instance & Chapter 3, Chapter 2, Section 2 \\
\(\exists x, P\, x\) (existential, \texttt{Prop}-valued) & \texttt{∃ x, P x}, \emph{not} Σ, a proposition witnessed by a pair, but with no witness-extraction since it lives in \texttt{Prop} & Chapter 4 \\
proof-irrelevant universe & \texttt{Prop} & Chapter 4, Chapter 2, Section 2 \\
inductive data (beyond bare CoC) & \texttt{inductive}, \texttt{structure} & Chapters 1, 3, 12 \\
Curry--Howard: propositions-as-types & \texttt{Prop}/proof terms & Chapter 4 \\
negation \(\neg P\) & \texttt{¬P} (defined as \texttt{P → False}) & Chapter 4 \\
progress + preservation & the kernel type-checker of Lean & Chapter 6, Section 3 \\
Church numerals (encoded data) & \texttt{Nat} (but built via \texttt{inductive}, not encoded) & Chapter 2, Section 1, Chapter 14 (aside) \\
Church booleans / if-then-else via application & \texttt{Bool}, \texttt{if}/\texttt{match} (but native, not encoded) & Chapter 14 (aside) \\
\end{longtable}

Two related facts, explained where their own dictionary rows live rather than repeated here. Tactics (\hyperref[sec:02-terminology-and-coc:01-terminology:1]{Chapter 2, Section 1}) add nothing to the underlying calculus; they are a user interface for building the same terms step by step. Elaboration/unification (\hyperref[sec:02-terminology-and-coc:01-terminology:1]{Chapter 2, Section 1}) is type inference for this calculus, a well-understood algorithm, not magic.
